% macros.tex — THE ONE PLACE MATH MACROS LIVE
% ---------------------------------------------------------------------------
% • PDF: this file is read as-is into the LaTeX preamble (cjgg-pdf format).
% • HTML: macros-to-mathjax.py converts it into MathJax config, so every
%   page has every macro with no per-page include. Run it after editing
%   (the homepage's pre-render hook does this automatically).
% Rules the converter understands: \providecommand / \newcommand /
% \renewcommand, with [n] arguments and one [default] optional argument.
% Nothing else (no \ifdefined, no \let, no \DeclareMathOperator).
% ---------------------------------------------------------------------------

% logic
\providecommand{\SuchThat}{\operatorname{s.t.}}

% set theory
\providecommand{\Image}{\operatorname{im}}
\providecommand{\Into}{\hookrightarrow}
\providecommand{\Onto}{\twoheadrightarrow}
\providecommand{\Isomorphism}{\xrightarrow{\cong}}
\providecommand{\Functions}{\operatorname{Func}}
\providecommand{\Indicator}[1]{\mathbb{1}_{#1}}

% blackboard bold
\providecommand{\Ab}{\mathbb{A}}
\providecommand{\Bb}{\mathbb{B}}
\providecommand{\Cb}{\mathbb{C}}
\providecommand{\Db}{\mathbb{D}}
\providecommand{\Eb}{\mathbb{E}}
\providecommand{\Fb}{\mathbb{F}}
\providecommand{\Gb}{\mathbb{G}}
\providecommand{\Hb}{\mathbb{H}}
\providecommand{\Ib}{\mathbb{I}}
\providecommand{\Jb}{\mathbb{J}}
\providecommand{\Kb}{\mathbb{K}}
\providecommand{\Lb}{\mathbb{L}}
\providecommand{\Mb}{\mathbb{M}}
\providecommand{\Nb}{\mathbb{N}}
\providecommand{\Ob}{\mathbb{O}}
\providecommand{\Pb}{\mathbb{P}}
\providecommand{\Qb}{\mathbb{Q}}
\providecommand{\Rb}{\mathbb{R}}
\providecommand{\Sb}{\mathbb{S}}
\providecommand{\Tb}{\mathbb{T}}
\providecommand{\Ub}{\mathbb{U}}
\providecommand{\Vb}{\mathbb{V}}
\providecommand{\Wb}{\mathbb{W}}
\providecommand{\Xb}{\mathbb{X}}
\providecommand{\Yb}{\mathbb{Y}}
\providecommand{\Zb}{\mathbb{Z}}

% group theory
\providecommand{\Pure}[1]{\mathcal{P}_{#1}}
\providecommand{\Braid}[1]{\mathcal{B}_{#1}}
\providecommand{\Symmetric}[1]{\mathcal{S}_{#1}}
\providecommand{\Cyclic}[1]{\mathcal{C}_{#1}}
\providecommand{\Modular}[1]{\mathbb{Z}/{#1}\mathbb{Z}}
\providecommand{\Dihedral}[1]{\mathcal{D}_{#1}}
\providecommand{\Alternating}[1]{\mathcal{A}_{#1}}
\providecommand{\Klein}{\mathcal{V}}
\providecommand{\Quaternion}[1][8]{\mathcal{Q}_{#1}}
\providecommand{\Affine}[2]{\operatorname{Aff}_{#1}({#2})}
\providecommand{\Permutation}[1]{\operatorname{Perm}(#1)}
\providecommand{\ConjugacyClass}[2][]{\operatorname{cl}_{#1}({#2})}
\providecommand{\Center}[1]{\mathbf{Z}(#1)}
\providecommand{\Normal}{\unlhd}
\providecommand{\Acts}{⟳}
\providecommand{\Fix}{\operatorname{Fix}}
\providecommand{\Stabilizer}[2][]{\operatorname{Stab}_{#1}({#2})}
\providecommand{\Orbit}[2][]{\operatorname{Orb}_{#1}({#2})}
\providecommand{\Centralizer}[2][]{\operatorname{C}_{#1}({#2})}
\providecommand{\Abelianization}[1]{{#1}^{\mathrm{Ab}}}
\providecommand{\Sign}{\operatorname{sgn}}

% algebra
\providecommand{\Hom}{\operatorname{Hom}}
\providecommand{\End}{\operatorname{End}}
\providecommand{\Aut}{\operatorname{Aut}}
\providecommand{\Identity}{\mathbb{1}}
\providecommand{\Id}{\mathbb{I}}

% fields
\providecommand{\Characteristic}{\operatorname{char}}

% linear algebra
\providecommand{\Mat}{\operatorname{Mat}}
\providecommand{\Span}{\operatorname{Span}}
\providecommand{\Null}{\operatorname{Null}}
\providecommand{\Rank}{\operatorname{rank}}
\providecommand{\Trace}{\operatorname{Tr}}
\providecommand{\Diagonal}{\operatorname{Diag}}
\providecommand{\GL}{\operatorname{GL}}
\providecommand{\SL}{\operatorname{SL}}
\providecommand{\Orthogonal}{\operatorname{O}}
\providecommand{\SO}{\operatorname{SO}}
\providecommand{\Unitary}{\operatorname{U}}
\providecommand{\SU}{\operatorname{SU}}
\providecommand{\PU}{\operatorname{PU}}
\providecommand{\PSL}{\operatorname{PSL}}
\providecommand{\PGL}{\operatorname{PGL}}
\providecommand{\PSU}{\operatorname{PSU}}

\providecommand{\SymmetricPower}[2][]{\operatorname{Sym}^{#1}({#2})}
\providecommand{\ExteriorPower}[2][]{\wedge^{#1}({#2})}

% fourier theory
\providecommand{\ClassFunctions}[1]{\mathscr{C}({#1})}
\providecommand{\ContinuousFunctions}[1][]{C^{#1}}
\providecommand{\IntegrableFunctions}[2][2]{L^{#1}({#2})}
\providecommand{\Fourier}{\mathscr{F}}

% algebraic geometry
\providecommand{\Disc}[1]{\mathrm{disc}({#1})}
\providecommand{\Res}[2]{\mathrm{Res}({#1}, {#2})}
\providecommand{\Gal}[1]{\mathrm{Gal}(#1)}
\providecommand{\rationalto}{\dashrightarrow}
\providecommand{\blowup}[2]{\mathrm{Bl}_{#1}({#2})}
\providecommand{\Grassmannian}[2]{\mathrm{Gr}_{#1}({#2})}
\providecommand{\Frob}[1]{\mathrm{Frob}_{#1}}
\providecommand{\Poly}[2]{\mathrm{Poly}_{#1}({#2})}

% category theory
\providecommand{\catname}[1]{\mathsf{#1}}
\providecommand{\VectorSpaces}[1]{\catname{Vect}_{#1}}

% resolvent degree
\providecommand{\RD}{\mathrm{RD}}

% complex analysis
\providecommand{\abs}[1]{\left|{#1}\right|}
\providecommand{\Arg}[1]{\mathrm{Arg}\left({#1}\right)}   % principal argument; lowercase \arg is LaTeX's own

% vectors (Math 120)
% Capitalised like the rest of the library: LaTeX already has a lowercase
% \vector (picture-mode arrows), which \providecommand would silently keep.
\providecommand{\Vector}[1]{\langle #1 \rangle}
\providecommand{\norm}[1]{\left\lVert #1 \right\rVert}
\providecommand{\iv}{\mathbf{i}}
\providecommand{\jv}{\mathbf{j}}
\providecommand{\kv}{\mathbf{k}}
\providecommand{\comp}[2]{\operatorname{comp}_{#1}\left(#2\right)}
\providecommand{\proj}[2]{\operatorname{proj}_{#1}\left(#2\right)}

% calculus (diffcoeff defines \diff/\diffp in PDF; these are the HTML fallbacks.
% The cjgg-pdf format loads diffcoeff BEFORE this file so \providecommand defers to it.)
\providecommand{\dl}[1]{\mathop{}\!\mathrm{d}{#1}}
\providecommand{\dln}[2]{\mathop{}\!\mathrm{d}^{#2}{#1}}
\providecommand{\diff}[2]{\frac{\dl{#1}}{\dl{#2}}}
\providecommand{\diffn}[3]{\frac{\dln{#1}{#3}}{\dl{#2}^{#3}}}
\providecommand{\dlp}[1]{\mathop{}\!\partial {#1}}
\providecommand{\diffp}[2]{\frac{\dlp{#1}}{\dlp{#2}}}
\providecommand{\diffpn}[3]{\frac{\partial^{#3} {#1}}{\partial #2^{#3}}}
\providecommand{\Divergence}{\operatorname{div}}
\providecommand{\Curl}{\operatorname{curl}}
\providecommand{\Gradient}{\operatorname{grad}}
\providecommand{\D}{\mathrm{D}}
\providecommand{\I}{\mathrm{I}}
